% titel: Blatt 10, A4: Gleichgewichtslage bestimmen
% Teil II: der durchdachte Weg.

\subsection*{a) HSO$_4^-$ + CO$_3^{2-} \rightleftharpoons SO_4^{2-} + HCO_3^-$}
Paar: HSO$_4^-$/SO$_4^{2-}$ ($pK_{a2}\approx1.9$), HCO$_3^-$/CO$_3^{2-}$
($pK_a\approx10.3$). HSO$_4^-$ ist die stärkere Säure, CO$_3^{2-}$ die
stärkere Base. Das Proton wandert von HSO$_4^-$ zu CO$_3^{2-}$
$\Rightarrow$ Gleichgewicht stark auf der \textbf{rechten} Seite.

\subsection*{b) HPO$_4^{2-}$ + H$_2$O $\rightleftharpoons$ H$_2$PO$_4^-$ + OH$^-$}
HPO$_4^{2-}$ ($pK_a\approx12.4$ als Säure) ist eine sehr schwache Säure,
H$_2$O ($pK_a\approx14$... H$_2$PO$_4^-$ $pK_a\approx7.2$) dazwischen. Das
Proton von HPO$_4^{2-}$ an H$_2$O zu geben erzeugt H$_2$PO$_4^-$ + OH$^-$,
aber H$_2$O ist eine zu schwache Base. $pK_a(\mathrm{HPO_4^{2-}}-
\mathrm{H_2PO_4^-}) \approx 7.2$ als Säure von H$_2$PO$_4^-$ --- so gilt:
H$_2$PO$_4^-$ ist stärker sauer als H$_2$O. Gleichgewicht auf der Seite der
schwächsten Säure/Base $\Rightarrow$ \textbf{links}.

\subsection*{c) NH$_4^+$ + OH$^- \rightleftharpoons$ NH$_3$ + H$_2$O}
Paare: NH$_4^+$/NH$_3$ ($pK_a\approx9.3$), H$_2$O/OH$^-$ ($pK_a(\mathrm{H_2O})
\approx14$, $pK_a(\mathrm{H_3O^+})\approx-1.7$... vergleich über Basenstärke:
OH$^-$ ist eine viel stärkere Base als H$_2$O). NH$_4^+$ reagiert mit dem
starken OH$^-$ zu NH$_3$ + H$_2$O. Das Proton geht auf die stärkere Base
(OH$^-$) $\Rightarrow$ Gleichgewicht deutlich auf der
\textbf{rechten} Seite.

\begin{ergebnis}
a) rechts (HSO$_4^-$ + CO$_3^{2-}$ $\to$ SO$_4^{2-}$ + HCO$_3^-$);
b) links (HPO$_4^{2-}$ + H$_2$O $\to$ H$_2$PO$_4^-$ + OH$^-$, schwach);
c) rechts (NH$_4^+$ + OH$^- \to$ NH$_3$ + H$_2$O).
\end{ergebnis}

\begin{kontrolle}
Die beiden starken Basen CO$_3^{2-}$ (a) und OH$^-$ (c) ziehen das Proton
an, Gleichgewicht rechts. In (b) konkurrieren zwei schwache Systeme, die
ungünstigere Protonierung (H$_2$PO$_4^-$ als Säure) bestimmt die linke
Seite. Konsistent mit den $pK_a$.
\end{kontrolle}

\begin{merksatz}
Ein Proton wandert zur stärkeren Base; Gleichgewicht liegt auf der Seite
der schwächeren Säure und schwächeren Base. $pK_a$-Vergleich entscheidet.
\end{merksatz}
